\section{Conclusion}\label{sec:conclusion}

Cartel detection has for decades relied on signals extracted from the bids themselves --- variance screens, rotation patterns, bid-distribution shapes, non-parametric tests of competitive behavior. This paper shows that two structural features of the firms involved, orthogonal to anything in the bid stream, can carry substantial detection signal. \emph{Pair-level worker flow} measures the count of workers moving between firms across consecutive years in the matched employer--employee record; \emph{product-market overlap} measures the number of procurement items on which the two firms both bid. Each feature is grounded in a distinct theoretical primitive: the former in the labor-economics literature on inter-firm knowledge transmission through worker mobility \citep{stoyanov2012worker, jarosch2021learning, sorkin2018ranking}; the latter in the canonical cartel-theory observation \citep{stigler1964theory, green1984noncooperative} that cartels must coordinate over a shared product scope.

A class-balanced logistic composite of the two features achieves a leave-one-CADE-cartel-out macro AUC of $0.845$ on the six-cartel S\~{a}o Paulo validation sample, compared to $0.734$ for product-market overlap alone. The incremental contribution of worker flow is roughly $+0.11$ macro AUC points. The cluster bootstrap over cartels places a 95\% confidence interval on this increment of $[-0.07, +0.27]$ and an empirical $p$-value against the null of zero gain of $0.10$. A label-permutation test on the composite's pooled AUC rejects the no-signal null at the 2\% level. Taken together, the two inference tests say: the composite captures a real signal that bid-based screens miss, and the specific contribution of worker flow to that signal is economically meaningful but at the margin of what a five-cartel validation design can statistically resolve.

The paper is explicit about where the composite succeeds and where it does not. The screen generalizes strongly on four of the six CADE cartels (fold AUCs between $0.84$ and $1.00$), including the three cartels with the largest within-cartel pair counts. It generalizes weakly on the two shortest-lived cartels (aquecedores solares, 2009--2010; cafeteria aeroporto, 2014), yielding fold AUCs of approximately $0.66$. This failure mode is consistent with the framework's intuition that structural fingerprints accumulate over time and that one-year cartels will not yet have left detectable traces in either the labor-mobility or the product-market-overlap dimension. It defines an operating envelope for the methodology: mature, multi-year cartels with consolidated product scope, rather than short-lived episodes of price fixing.

Three implications organize the paper's contribution.

\emph{First, the screen is a prototype tool, not a turnkey deployment system.} Computation of both features requires (i)~a matched employer--employee register linked to the procurement bidder universe and (ii)~a record of which firms bid on which items. No bid values, reference prices, or auction-format details are needed. Competition authorities with access to these two data layers can compute the features on their own universe in a single pass. However, producing a ranked audit priority score requires calibration on ground-truth cartel data specific to the jurisdiction, and the classifier coefficients will not transfer without re-estimation. Our six-cartel calibration is a first demonstration of feasibility, not a universally applicable scoring rule.

\emph{Second, the two fingerprints fill an operating gap left by bid-based screens in dynamic-auction settings.} We show in Section~\ref{sec:rob_variance} that the classical variance-based cartel screen of \citet{abrantesmetz2006screening} produces a pair-level AUC below chance on BEC data. We attribute this failure to a structural feature of the preg\~{a}o eletr\^{o}nico: the descending-auction mechanism compresses within-auction bid distributions across cartel and non-cartel auctions alike, eliminating the variance signal that makes the Abrantes-Metz screen informative in sealed-bid settings. Jurisdictions that have adopted dynamic electronic procurement auctions --- which is an increasing share of public procurement in Latin America, Asia, and multilateral agencies --- cannot rely on variance-based screens, and the firm-pair fingerprints we propose are, to our knowledge, the first applicable alternative in this setting.

\emph{Third, the principal constraint on inference is the small number of convicted cartels in our ground truth.} With five effective cartels in the validation set, cluster-bootstrap confidence intervals on the composite's macro AUC span roughly $0.3$ AUC points. Expanding the ground truth --- by incorporating CADE cases beyond the SP-relevant subset, by linking to other national competition authorities that maintain conviction registers, or by pooling across multiple Latin American jurisdictions --- would sharpen both the point estimates and the characterization of failure modes. The infrastructure we build is directly reusable in any country that has both a public procurement platform recording bidder identities and a matched employer--employee register. To our knowledge we are the first to implement both-sides-linked cartel detection in a developing-country setting, and we hope the methodology will be taken up and extended by others.

\subsection*{What our inference can and cannot resolve}

We close with a candid statement about the statistical evidence.

The composite's pooled AUC of $0.809$ is well above what a label-permutation null produces ($p = 0.020$). The composite's macro AUC of $0.845$ exceeds every reference specification we tested: overlap alone, static labor cosine, CNAE-match-plus-overlap, and common-auction-plus-overlap. The raw-scale pair-level worker flow differs by a factor of $24{,}000$ between cartel and non-cartel pairs ($74$ workers versus $0.003$ workers on average). These three pieces of evidence, taken together, make a strong case that convicted cartels leave a detectable fingerprint in pair-level labor mobility and product-market participation.

The specific test of whether the worker-flow feature contributes to the composite \emph{beyond} what product-market overlap alone already captures is harder. The cluster bootstrap over cartels gives a 95\% CI on that increment of $[-0.07, +0.27]$ with empirical $p \approx 0.10$. A reader who applies a strict 5\% bootstrap-difference threshold will conclude that the incremental contribution of worker flow is not statistically distinguishable from zero on this sample. A reader who weighs the economic magnitude ($+0.11$ macro AUC points), the consistent direction of the effect across the bootstrap distribution, the strong permutation-test result on the composite as a whole, and the conceptual alignment with the labor-economics literature on inter-firm knowledge transmission will reach a different conclusion. Our own reading tilts toward the latter, but the evidence is not strong enough to force that reading on others.

What is unambiguous, and what we view as the paper's most robust finding, is that structural firm-pair fingerprints detect convicted cartels in a setting where bid-based screens do not work. The precision limits on the specific labor contribution will sharpen with larger ground-truth samples. The qualitative finding --- that cartel detection outside the bid stream is feasible on preg\~{a}o data --- stands.

\subsection*{Selection caveats}

Two selection caveats apply to any reading of our numbers, and we state them explicitly.

The ground truth is a set of \emph{caught} cartels, not a random sample of all cartels. Unconvicted cartels sit in our negative class, adding noise and biasing the measured AUC downward. Our point estimates should therefore be read as lower bounds on what the composite could achieve if the ground truth were complete. Conversely, the six CADE convictions we use are the result of an enforcement process that selects for cartels detectable by existing methods, and our alternative fingerprints may be overfit to that selection --- meaning they may be less informative for cartels of a different structure that go uncaught. The honest reading of our $0.84$ macro AUC is therefore: \emph{this is what pair-level worker flow plus product-market overlap can achieve on the cartels that the existing enforcement infrastructure has already found}. Whether they can find cartels that the existing infrastructure has missed is a question we cannot answer without a field experiment, and one we leave for future research.

The second caveat concerns the calibration of the composite across jurisdictions. Our logistic coefficients are estimated on Brazilian data and embedded in the feature-standardization scale used at fit time. A competition authority in Mexico, Colombia, or the European Union that wants to deploy the screen should not import our coefficients directly; they should re-estimate on their own ground truth and re-standardize their features. The architecture of the screen --- two features, class-balanced logistic, leave-one-cartel-out calibration --- transfers across jurisdictions. The specific numeric coefficients do not.

Despite these caveats, the paper's core message is that structural firm-pair data carry cartel-detection signal, that this signal is orthogonal to the bid-distribution signals that existing screens exploit, and that in dynamic-auction settings where bid-distribution screens are structurally ineffective, firm-pair fingerprints are a practical alternative. We think this is a useful addition to the applied detection toolkit and we look forward to seeing it extended to larger and more diverse ground-truth samples.
